\documentclass[12pt]{article}
\usepackage[T1]{fontenc}
\usepackage[utf8]{inputenc}
\usepackage{lmodern}
\usepackage{textcomp}
\usepackage{enumitem}
\usepackage{geometry}
\geometry{margin=1in}
\begin{document}
\begin{center}
Answer Key - Mathematics - Undergraduate Level Assessment 21 \
\small (Answers are ordered exactly as in the question paper.)
\end{center}

\section*{Multiple Choice}
\begin{enumerate}[label=\arabic*.]
\item \textbf{Answer:} C
\\ \textit{Options:} A) greater than 0.50; B) between 0.16 and 0.50; C) between 0.02 and 0.16; D) between 0.01 and 0.02
\\ \textit{Note:} The correct answer is between 0.02 and 0.16 because when a fair die is tossed 360 times, the number of sixes follows a binomial distribution, and applying the normal approximation to the binomial distribution reveals that the probability of obtaining 70 or more sixes is relatively low, falling within that range.
\item \textbf{Answer:} D
\\ \textit{Options:} A) $x^2$ + $y^2$ + $z^2$ + 8 = 6 y; B) ($x^2$ + $y^2$ + $z^2$$)^2$ = 8 + 36($x^2$ + $z^2$); C) ($x^2$ + $y^2$ + $z^2$ + $8)^2$ = 36($x^2$ + $z^2$); D) ($x^2$ + $y^2$ + $z^2$ + $8)^2$ = 36($x^2$ + $y^2$)
\\ \textit{Note:} The equation of the torus is derived from revolving the circle defined by (y - $3)^2$ + $z^2$ = 1 around the z-axis, resulting in the equation ($x^2$ + $y^2$ + $z^2$ + $8)^2$ = 36($x^2$ + $y^2$), which represents the surface of a torus with a specific radius and distance from the axis of rotation.
\item \textbf{Answer:} D
\\ \textit{Options:} A) True, True; B) True, False; C) False, True; D) False, False
\\ \textit{Note:} The correct answer is False for both statements because a local maximum can occur at a point where the derivative does not exist (such as at a cusp), and there are no non-constant continuous functions from R to Q due to the density of R in Q, which contradicts the intermediate value theorem.
\item \textbf{Answer:} B
\\ \textit{Options:} A) -2; B) -1; C) 0; D) 2
\\ \textit{Note:} The determinant of matrix A, which has entries defined by i + j, is -1, and the determinant of matrix B, a 3 by 3 matrix with the same entry definition, is 0, leading to a sum of -1 for detA + detB.
\item \textbf{Answer:} B
\\ \textit{Options:} A) True, True; B) False, False; C) True, False; D) False, True
\\ \textit{Note:} The correct answer is False for both statements because while R can be a splitting field of a polynomial over Q, it does not imply the existence of a field with 60 elements, as finite fields exist only for orders that are powers of prime numbers, and 60 is not a prime power.
\end{enumerate}

\section*{True / False}
\begin{enumerate}[label=\arabic*.]
\setcounter{enumi}{5}
\item \textbf{Answer:} False
\\ \textit{Options:} A) True; B) False
\\ \textit{Note:} The statement is false because the number of trailing zeros in a factorial is determined by the number of times 5 is a factor in the integers up to k, and multiple values of k can yield the same number of trailing zeros, including 99.
\item \textbf{Answer:} False
\\ \textit{Options:} A) True; B) False
\\ \textit{Note:} The probability that one integer will be the square of the other is 0, as there are infinitely many integers and only a limited number of pairs of integers that satisfy this condition, making the probability effectively zero.
\item \textbf{Answer:} True
\\ \textit{Options:} A) True; B) False
\\ \textit{Note:} The polynomial 4 x - 2 is irreducible over the field of rational numbers Q because it is a linear polynomial with no non-trivial rational roots other than its scalar multiple.
\item \textbf{Answer:} True
\\ \textit{Options:} A) True; B) False
\\ \textit{Note:} The order of the cyclic subgroup generated by 18 in Z\_24 is 4 because the smallest positive integer k such that 18 k is congruent to 0 modulo 24 is 4, as 18(4) = 72, which is equivalent to 0 in Z\_24.
\item \textbf{Answer:} False
\\ \textit{Options:} A) True; B) False
\\ \textit{Note:} The statement is false because for Z\_3[x]/($x^3$ + $x^2$ + c) to be a field, the polynomial $x^3$ + $x^2$ + c must be irreducible over Z\_3, and with c = 1, the polynomial factors as (x + 1)($x^2$ + 1), thus not being irreducible.
\end{enumerate}

\section*{Long Form Response}
\begin{enumerate}[label=\arabic*.]
\setcounter{enumi}{10}
\item \textbf{Answer:} Because the rocking chairs are indistinguishable from each other and the stools are indistinguishable from each other, we can think of first placing the two stools somewhere in the ten slots and then filling the rest with rocking chairs. The first stool has $10$ slots in which it can go, and the second has $9$. However, because they are indistinguishable from each other, we have overcounted the number of ways to place the stools by a factor of $2$, so we divide by $2$. Thus, there are $\frac{10 \cdot 9}{2} = \boxed{45}$ distinct ways to arrange the ten chairs and stools for a meeting.
\\ \textit{Note:} correct because it correctly applies the combinatorial principle of counting arrangements by first selecting positions for the indistinguishable stools among the total slots and then adjusting for overcounting due to the indistinguishability of the stools.
\item \textbf{Answer:} The implications of Statement 1, which asserts that every group of order \( p^2 \) (where \( p \) is prime) is Abelian, highlight a fundamental property of finite groups. Specifically, groups of such order can only be cyclic or direct products of cyclic groups, both of which are Abelian. This characteristic simplifies the structure and analysis of these groups in group theory. In relation to Statement 2, the conditions for a Sylow \( p \)-subgroup to be normal in a group \( G \) are crucial for understanding the group's structure. If the number of Sylow \( p \)-subgroups is one, then that subgroup is normal, which illustrates the relationship between subgroup structure and group properties. Both statements are true and underscore the importance of prime orders in the classification and behavior of groups.
\\ \textit{Note:} correct because it accurately identifies that groups of order \( p^2 \) are guaranteed to be Abelian, which simplifies their structure, and it highlights the importance of understanding the conditions under which Sylow \( p \)-subgroups are normal for analyzing the overall structure of a group \( G \).
\end{enumerate}

\vfill
\noindent\textit{End of Answer Key}
\end{document}